\documentclass[12pt]{article}
\usepackage{amsfonts,amsthm,amsmath,amssymb,mathtools}
\usepackage{mathrsfs}
\usepackage{enumitem}
\usepackage{tikz-cd}
\usepackage{geometry}
\usepackage{mlmodern}

\geometry{letterpaper, portrait, margin=1in}

\setcounter{section}{0}

\renewcommand{\thesection}{\S\arabic{section}}
\renewcommand{\thesubsection}{\arabic{section}}

\newtheorem{thm}{Theorem}
\newenvironment{theorem}[1]{
	\renewcommand{\thethm}{#1}
	\begin{thm}
		}{\end{thm}}

\newtheorem{lma}{Lemma}
\newenvironment{lemma}[1]{
	\renewcommand{\thelma}{#1}
	\begin{lma}
		}{\end{lma}}

\theoremstyle{definition}
\newtheorem{ex}{}
\newenvironment{exercise}[1]{
	\renewcommand{\theex}{#1}
	\begin{ex}
		}{\end{ex}}

\title{Homework Assignment I}
\author{Connor Mooneyhan}
\date{August 29, 2025}

\begin{document}

\maketitle

\begin{ex}
	Suppose that $f: [a, b] \to \mathbb{R}$ is a bounded function such that \begin{equation*}
		L(f, P, [a, b]) = U(f, P, [a, b])
	\end{equation*}
	for some partition $P$ of $[a, b]$.
	Prove that $f$ is a constant function on $[a, b]$.
\end{ex}
\begin{proof}
	Denote the hypothesized partition $P$ as $a = x_0, x_1, \dots, x_n = b$.
	We proceed to show that $f$ is constant on each subinterval $[x_{i-1}, x_i]$, where $i = 1, \dots, n$.

	Suppose that for some such $i$, there exist two points $p$ and $q$ in $[x_{i-1}, x_i]$ such that $f(p) < f(q)$.
	Then $\inf_{[x_{i-1}, x_i]} f \leq f(p) < f(q) \leq \sup_{[x_{i-1}, x_i]} f$.
	Since $\inf_{[x_{j-1}, x_j]} f \leq \sup_{[x_{j-1}, x_j]} f$ for any $j$, it follows that \begin{equation*}
		L(f, P, [a, b]) = \sum_{[x_{j-1}, x_j]} (x_j - x_{j-1})\inf_{[x_{j-1}, x_j]} f < \sum_{[x_{j-1}, x_j]} (x_j - x_{j-1})\sup_{[x_{j-1}, x_j]} f = U(f, P, [a, b]).
	\end{equation*}
	To drive the point home, the strict inequality comes from the fact that comparing the terms of the sum for a general $[x_{j-1}, x_j]$ yields an inclusive inequality, whereas for at least one pair of terms, namely the terms corresponding to $[x_{i - 1}, x_i]$, the inequality is strict.
	This inequality contradicts the hypothesis that $L(f, P, [a, b]) = U(f, P, [a, b])$, so we conclude that $f$ is constant on each subinterval $[x_{j - 1}, x_j]$.

	Since $f$ is constant on each subinterval, we have in particular that $f(x_{j - 1}) = f(x_j)$ for each $j$. This means that, in fact, $f(x_0) = f(x_1) = \cdots = f(x_n)$, so each subinterval is constant \textit{at the same value}. Thus, $f$ is constant on $[a, b]$.
\end{proof}

\begin{ex}
	Suppose $a \leq s < t \leq b$.
	Define $f: [a, b] \to \mathbb{R}$ by \begin{equation*}
		f(x) = \begin{cases}
			1 & \text{if } s < x < t, \\
			0 & \text{otherwise.}
		\end{cases}
	\end{equation*}
	Prove that $f$ is Riemann integrable on $[a, b]$ and that $\int_a^b f = t - s$.
\end{ex}
\begin{proof}
	We will define a sequence $(P_n)$ of partitions such that \begin{equation*}
		\lim_{n \to \infty}L(f, P_n, [a, b]) = \lim_{n \to \infty} U(f, P_n, [a, b]) = t - s.
	\end{equation*}
	Since \begin{equation*}
		\lim_{n \to \infty}L(f, P_n, [a,b]) \leq L(f, [a, b]) \leq U(f, [a, b]) \leq \lim_{n \to \infty}U(f, P_n, [a, b]),
	\end{equation*}
	this will squeeze the upper and lower Riemann integrals so that the above inequalities become equalities, allowing us to conclude that $f$ is Riemann integrable over $[a, b]$ and $\int_a^b f = t - s$.

	Let $P_0 = a, s, t, b$ and $P_1 = a, s, s + (t-s)/3, t - (t-s)/3, t, b$.
	Continuing in this fashion for all other $n$, letting $\varepsilon = (t-s)/3$, we define \begin{equation*}
		P_n = a, s, s + \frac{\varepsilon}{n}, s + \frac{\varepsilon}{n-1}, \dots, s + \varepsilon, t - \varepsilon, \dots, t - \frac{\varepsilon}{n-1}, t - \frac{\varepsilon}{n}, t, b.
	\end{equation*}
	For the ``outer'' intervals $[a, s]$ and $[t, b]$, the corresponding terms of the lower and upper Riemann sums agree since \begin{equation*}
		\inf_{[a, s]} f = \sup_{[a, s]} f = 0 = \inf_{[t, b]}f = \sup_{[t, b]}f.
	\end{equation*}
	The two ``limiting'' intervals $[s, s + \varepsilon/n]$ and $[t - \varepsilon/n, t]$ have the supremum being 1 and the infimum being 0.
	All the remaining ``inner'' intervals are strictly contained between $s$ and $t$, so the infimum and supremum are both 1.
	We will use this together with the fact that the lengths of all of these ``inner'' intervals add up to $t - \varepsilon/n - s - \varepsilon/n = t - s - 2\varepsilon/n$.
	Taking all of this into account, we get for $n > 1$ that \begin{align*}
		L(f, P_n, [a, b]) & = (s - a) \cdot 0 + \left(s + \frac{\varepsilon}{n} - s\right) \cdot 0 + (t - s - \frac{2\varepsilon}{n}) + \left( t - t + \frac{\varepsilon}{n} \right)\cdot 0 + (b - t) \cdot 0 \\
		                  & = t - s - \frac{2\varepsilon}{n}
	\end{align*}
	and \begin{align*}
		U(f, P_n, [a, b]) & = (s - a) \cdot 0 + \left(s + \frac{\varepsilon}{n} - s\right) \cdot 1 + (t - s - \frac{2\varepsilon}{n}) + \left( t - t + \frac{\varepsilon}{n} \right)\cdot 1 + (b - t) \cdot 0 \\
		                  & = \frac{\varepsilon}{n} + t - s - \frac{2\varepsilon}{n} + \frac{\varepsilon}{n}                                                                                                  \\
		                  & = \frac{\varepsilon}{n} + t - s - \frac{2\varepsilon}{n} + \frac{\varepsilon}{n}                                                                                                  \\
		                  & = t - s.
	\end{align*}
	As $n$ goes to infinity, both the lower and upper sums limit to $t - s$, as desired.
\end{proof}

\begin{ex}
	Suppose $f: [a, b] \to \mathbb{R}$ is a bounded function.
	Prove that $f$ is Riemann integrable if and only if for each $\varepsilon > 0$ there exists a partition $P$ of $[a, b]$ such that \begin{equation} \label{3.1}
		U(f, P, [a, b]) - L(f, P, [a, b]) < \varepsilon.
	\end{equation}
\end{ex}
\begin{proof}
	Taking the forward implication first, suppose that for some $\varepsilon > 0$, there is no partition $P$ of $[a, b]$ which satisfies \eqref{3.1}.
	Then $U(f, P, [a, b]) - L(f, P, [a, b]) \geq \varepsilon$ for all partitions $P$.
	In the limit (over all $P$), we get $U(f, [a, b]) - L(f, [a, b]) \geq  \varepsilon$, so in particular $U(f, [a, b]) \neq L(f, [a, b])$.

	In the other direction, given $\varepsilon > 0$, choose a partition $P$ such that \begin{equation*}
		U(f, [a, b]) \leq U(f, P, [a, b]) < L(f, P, [a, b]) + \varepsilon \leq L(f, [a, b]) + \varepsilon,
	\end{equation*}
	where the second inequality comes from \eqref{3.1}.
	Since $\varepsilon$ was arbitrary, we conclude that $U(f, [a, b]) \leq L(f, [a, b])$, which when combined with the standard inequality in the other direction, yields the desired equality.
\end{proof}

\begin{ex}
	Suppose $f, g: [a, b] \to \mathbb{R}$ are Riemann integrable.
	Prove that $f + g$ is Riemann integrable on $[a, b]$ and \begin{equation*}
		\int_a^b (f + g) = \int_a^b f + \int_a^b g.
	\end{equation*}
\end{ex}
\begin{proof}
	Fix $\varepsilon > 0$, and choose partitions $P^{\prime}$ and $P^{\prime\prime}$ such that \begin{equation*}
		U(f, P^{\prime}, [a, b]) - L(f, P^{\prime}, [a, b]) < \varepsilon/2
	\end{equation*}
	and \begin{equation*}
		U(g, P^{\prime\prime}, [a, b]) - L(g, P^{\prime\prime}, [a, b]) < \varepsilon/2.
	\end{equation*}
	Then the above inequalities also hold for $P = P^\prime \cup P^{\prime\prime}$, and we can add them to get \begin{equation} \label{4.1}
		U(f, P, [a, b]) + U(g, P, [a, b]) - \left( L(f, P, [a, b]) + L(g, P, [a, b]) \right) < \varepsilon.
	\end{equation}
	Write $P$ as $a = x_0, x_1, \dots, x_n = b$.
	Then for each $i = 1, \dots, n$, we have \begin{equation*}
		\left( \inf_{[x_{i - 1}, x_i]} f \right) + \left( \inf_{[x_{i - 1}, x_i]} g \right) \leq \inf_{[x_{i-1}, x_i]} \left( f + g \right),
	\end{equation*}
	so \begin{equation} \label{4.2}
		L(f, P, [a, b]) + L(g, P, [a, b]) \leq L(f + g, P, [a, b]).
	\end{equation}
	A similar argument shows that \begin{equation} \label{4.3}
		U(f, P, [a, b]) + U(g, P, [a, b]) \geq U(f + g, P, [a, b]).
	\end{equation}
	Combining these with \eqref{4.1}, we conclude that \begin{equation*}
		U(f + g , P, [a, b]) - L(f + g, P, [a, b]) < \varepsilon.
	\end{equation*}
	By problem 3, then, since $\varepsilon$ is arbitrary, $f + g$ is Riemann integrable.

	We now extend the construction of $P$ above to a sequence $(P_n)$, where each $P_n$ is a partition satisfying \begin{equation*}
		U(f + g, P_n, [a, b]) - L(f + g, P_n, [a, b]) < \varepsilon/n,
	\end{equation*}
	which gives us \begin{equation*}
		\lim_{n \to \infty} U(f + g, P_n, [a, b]) = U(f + g, [a, b])
	\end{equation*} and \begin{equation*}
		\lim_{n \to \infty} L(f + g, P_n, [a, b]) = L(f + g, [a, b]).
	\end{equation*}
	We now combine \eqref{4.2} and \eqref{4.3} as follows, adopting a more general partition from our new $(P_n)$ sequence: \begin{align*}
		L(f, P_n, [a, b]) + L(g, P_n, [a, b]) & \leq L(f + g, P_n, [a, b])                  \\
		                                      & \leq L(f + g, [a, b])                       \\
		                                      & \leq U(f + g, [a, b])                       \\
		                                      & \leq U(f + g, P_n, [a, b])                  \\
		                                      & \leq U(f, P_n, [a, b]) + U(g, P_n, [a, b]).
	\end{align*}
	Taking limits as $n \to \infty$ and removing redundant steps, we get \begin{align*}
		L(f, [a, b]) + L(g, [a, b]) & \leq L(f + g, [a, b])             \\
		                            & \leq U(f + g, [a, b])             \\
		                            & \leq U(f, [a, b]) + U(g, [a, b]).
	\end{align*}
	Since the first and last expressions are equal, all the inequalities become equalities, and we conclude that \begin{equation*}
		\int_a^b f + \int_a^b g = \int_a^b (f + g).
	\end{equation*}
\end{proof}

\begin{ex}
	The Banach-Tarski paradox refers to a counter-intuitive theorem that Stefan Banach and Alfred Tarski proved in 1924.
	It shows that a solid ball can be split up into 5 pieces which can then be reconstituted to form two balls which are both identical to the first one.
	This betrays our geometric intuition because in general, we want the union of two geometric objects to have volume equal to the sum of the volumes of the two objects.
	In this case, this principle does not hold since each of the duplicate balls has the same volume as the original one, but taking their union in a clever way results in that very same volume, which is one half what the volume of the union ``should'' be.
\end{ex}

\end{document}
